\section{Guidelines for Human-AI Interaction}
\paperhead{Guidelines for Human-AI Interaction}{Saleema Amershi, Dan Weld, Mihaela Vorvoreanu, Adam Fourney, Besmira Nushi, et al.}{CHI 2019 (ACM SIGCHI Conference on Human Factors in Computing Systems)}{https://doi.org/10.1145/3290605.3300233}{Crossref DOI 10.1145/3290605.3300233, title match 1.00, cited\_by=1886}

\subsection{What the paper does}

AI-infused systems act under uncertainty, produce false positives and false negatives, and change their behaviour as they learn, so they routinely violate classical usability heuristics such as consistency and error prevention. Two decades of design guidance for such systems exists, but the authors argue it is effectively unusable: it is scattered across AAAI, UbiComp, RecSys, SIGIR, HRI and KDD, and is often buried in the discussion or future-work section of a paper rather than stated as a guideline a designer could look up.

The contribution is a codification of more than 150 AI-related design recommendations into 18 short, inspectable guidelines, plus a systematic validation of that set. Each guideline is deliberately written as a rule of action of roughly 3 to 10 words beginning with a verb, carries a one-sentence qualifier, and contains no conjunctions, so an evaluator can decide by looking at an interface whether it was applied or violated. The 18 are grouped by when they bear on the interaction: \emph{Initially} (G1 make clear what the system can do, G2 make clear how well it can do it), \emph{During interaction} (G3 time services based on context, G4 show contextually relevant information, G5 match relevant social norms, G6 mitigate social biases), \emph{When wrong} (G7 support efficient invocation, G8 support efficient dismissal, G9 support efficient correction, G10 scope services when in doubt, G11 make clear why the system did what it did), and \emph{Over time} (G12 remember recent interactions, G13 learn from user behavior, G14 update and adapt cautiously, G15 encourage granular feedback, G16 convey the consequences of user actions, G17 provide global controls, G18 notify users about changes). Broad principles such as ``build trust'' were excluded precisely because they cannot be checked by inspection.

\subsection{How the method works}

The set was produced by a four-phase process. \textbf{Phase 1 (consolidation)} gathered recommendations from three sources: an audit of industry AI products and internally and externally asserted guidelines cross-referenced with customer feedback, public articles and editorials, and scholarly work. This yielded 168 candidate guidelines; three team members ran asynchronous affinity diagramming to cluster them into 35 concepts, then filtered out concepts judged too vague to design for, too specific to one AI scenario, or not AI specific, leaving 20 summarised in a sentence each.

\textbf{Phase 2 (modified heuristic evaluation)} adapted Nielsen-style heuristic evaluation so that the object under test was the guideline set rather than an interface: eleven team members each chose an AI-infused product or feature and spent one hour recording both applications and violations of the 20 guidelines, across 13 products in total, while reflecting on the guidelines themselves. Guidelines with no observable instances were dropped, redundant ones were merged, and conjunctive ones were split, for example a single ``allow efficient invocation, correction, and dismissal'' became G7, G8 and G9. The result was 18.

\textbf{Phase 3 (user study)} tested those 18 with 49 HCI practitioners. Each was assigned an AI-driven feature of a product they were familiar with and, for every guideline, first decided whether it ``does not apply'' to that feature and otherwise supplied a concrete example plus a rating on a 5-point semantic differential from ``clearly violated'' to ``clearly applied''; a closing questionnaire rated each guideline's wording from ``very confusing'' to ``very clear''. The products were 20 popular AI-infused applications selected by maximum-variance sampling into 10 categories with 2 per category, and 2 to 3 evaluators were assigned per product. Responses were then cleaned in a two-pass review by two team members: 55 duplicate instances, 73 misuses of ``does not apply'' and 14 disputed cases were removed, and 20 further instances were reclassified as violations. Analysis used four constructs: total applications plus violations as evidence of relevance; a large positive application-minus-violation gap as evidence of a guideline being widely implemented; a large negative gap as evidence of an important but neglected guideline; and high ``does not apply'' counts as evidence of scope limits. \textbf{Phase 4} had 11 experts compare old against revised wording for the 9 guidelines rewritten after Phase 3.

\subsection{What it reports}

Participants found at least one application or violation of every guideline in every product category, which the authors take as broad evidence of relevance across the tested spectrum. After cleaning, the study yielded 785 examples: 313 applications, 277 violations, 89 neutral ratings and 106 ``does not apply''. G11, on explaining why the system behaved as it did, drew one of the highest violation counts together with one of the lowest ``does not apply'' counts, that is, users could imagine an explanation almost everywhere and rarely got one. G3, G5 and G6 attracted the most ``does not apply'', and G5 and G6 were rated least clear, with participants disagreeing about whether social norms and bias were even in scope for the same product category. Some guidelines proved hard to observe in a single session without knowledge of the underlying algorithm: no instances of G10 were reported for either social network tested and none of G14 for either activity tracker. In Phase 4, experts preferred the revised wording for every revised guideline except G15, though five of the eleven still found G1 and G2 hard to tell apart. The authors note that 8 of the 18 guidelines map onto Horvitz's mixed-initiative principles.

\subsection{How Mytheus uses this}
\begin{takeaway}
Use the 18 guidelines as a second, qualitative evaluation instrument beside the golden-set CI harness: run a structured expert walkthrough of each of the six agents' surfaces with 2 to 3 evaluators per surface, scoring on the same 5-point clearly-violated to clearly-applied scale, and report guideline-level violation counts per agent. Several guidelines are citable design rationale for choices already frozen in the architecture: G10 (scope services when in doubt) is exactly the 0.4 to 0.7 silent-log band and the ContextGapAgent's silence terminal in the fallback chain; G8 and G9 (efficient dismissal and correction) are the acceptance test for the NO SILENT COMMITS candidate-then-confirm contract, so dismissing a candidate must cost one action and never a modal; G2 (make clear how well the system can do what it can do) means the 0.7 gate must surface calibrated confidence to the user, which the Brier score in the eval harness already measures; G11 is the citation contract, so a ContextGapAgent fill must arrive with its resolvable provenance attached; G18 covers notifying tenants when prompt versions ship through the 48-hour canary. The paper also warns about a confusion Mytheus can pre-empt: G15 and G17 were repeatedly conflated, so per-item feedback on a candidate and a global per-user control over which agents may surface at all must be visibly separate controls. Be clear about the limits: these are inspection heuristics, not runtime mechanisms. G3 gives no decision rule for when ProactiveAgent should interrupt versus defer, nothing in the set is evaluated against a response deadline so it offers nothing for the 1.5 second peer round trip budget, and G13 and G14 assume cross-session learning that in Mytheus must stay inside the org partition, a constraint the paper never considers.
\end{takeaway}

\subsection{Limitations and what it leaves open}
\begin{itemize}[leftmargin=*, itemsep=1pt, topsep=2pt]
  \item Everything is validated by static, single-session inspection of shipped consumer products, so the time-dependent guidelines are barely observable: participants reported no instances of G14 for either activity tracker and noted that G10 and G14 need a lengthy period with the product to assess, leaving the set untested for systems whose output evolves within a live session over a streaming transcript.
  \item The unit of analysis is one interface producing one shared output for a generic user; none of the 20 products personalises output per participant inside a shared session, so the guidelines give no account of how capability and confidence claims (G1, G2) should be phrased when two people in the same meeting legitimately receive different results.
  \item There is no latency, cost or throughput dimension anywhere: a guideline is scored only as applied or violated, never against a response deadline or a budget, so G3 and G10 supply no operational threshold for an interrupt-versus-defer policy under a p95 constraint.
  \item All 20 products are single-system interfaces, so cross-user and agent-to-agent information flow is entirely out of scope: the set says nothing about disclosing to a user that another participant's agent queried their agent, nor about consent or tenant-scoped authorisation for such a query.
  \item The validation establishes that the guidelines are recognisable and clearly worded, not that following them improves outcomes: there is no controlled comparison of guideline-conformant against non-conformant designs on task success, trust or error recovery, and no inter-rater reliability statistic is reported for the 785 judgements despite 56 recorded misinterpretations between guideline pairs.
\end{itemize}
